\begin{tikzpicture}[
  font=\scriptsize,
  kotak/.style={draw, rounded corners, minimum height=1.25cm, text width=2.0cm,
                align=center, font=\scriptsize\bfseries},
  panah/.style={-{Stealth[length=2mm]}, thick}
]
  \node[kotak, fill=blue!8] (aplikasi) at (0,0)
    {Aplikasi\\Medikara+\\{\tiny\mdseries di ponsel pasien}};
  \node[kotak, fill=orange!12] (api) at (6.2,0) {API jadwal\\dokter};
  \node[kotak, fill=blue!8] (simrs) at (9.6,0)
    {SIMRS\\{\tiny\mdseries Rumah Sakit Bekasi}};

  % Langkah 1 dan 4: percakapan antara aplikasi dan API.
  \draw[panah] ([yshift=0.3cm]aplikasi.east) --
    node[above, align=center, font=\scriptsize]
    {\textbf{1} request\\[1pt]
     {\ttfamily\tiny GET /jadwal-dokter}\\
     {\ttfamily\tiny ?rumah\_sakit=bekasi}\\
     {\ttfamily\tiny \&tanggal=2026-11-02}}
    ([yshift=0.3cm]api.west);
  \draw[panah] ([yshift=-0.3cm]api.west) --
    node[below, align=center, font=\scriptsize]
    {\textbf{4} response\\{\tiny kode status 200 OK dan isi JSON}}
    ([yshift=-0.3cm]aplikasi.east);

  % Langkah 2 dan 3: API mengambil data dari sistem di rumah sakit.
  \draw[panah] ([yshift=0.25cm]api.east) --
    node[above, font=\tiny] {\textbf{2} ambil} ([yshift=0.25cm]simrs.west);
  \draw[panah] ([yshift=-0.25cm]simrs.west) --
    node[below, font=\tiny] {\textbf{3} hasil} ([yshift=-0.25cm]api.east);

  % Isi response dalam format JSON.
  \node[draw, rounded corners, fill=green!10, align=left, font=\tiny\ttfamily,
        anchor=north, inner sep=5pt] (json) at (4.8,-1.45)
    {\{\\
     \quad "rumah\_sakit": "Rumah Sakit Medikara Bekasi",\\
     \quad "tanggal": "2026-11-02",\\
     \quad "jadwal": [\\
     \qquad \{"dokter": "dr. Rina", "poli": "Anak", "jam": "08:00-12:00",
       "sisa\_kuota": 6\},\\
     \qquad \{"dokter": "dr. Bayu", "poli": "Penyakit Dalam",
       "jam": "13:00-16:00", "sisa\_kuota": 0\}\\
     \quad ]\\
     \}};
  \node[font=\tiny, gray, anchor=south west] at (json.north west)
    {isi response yang dibaca aplikasi};
\end{tikzpicture}
